\section{DC-DFT Lagrangian과 전체 nuclear gradient}
\label{sec:dc-lagrangian}

\subsection{Constraint를 포함한 Lagrangian}

HF orbital coefficient를 $C$라 하고 HF stationarity residual의 VO block을
\begin{equation}
  f_{ai}^{\vk}(C,\vR)
  =\left[C_v^{\vk\dagger}F_{\mathrm{HF}}^{\vk}C_o^{\vk}\right]_{ai}
  \label{eq:hf-stationarity-residual}
\end{equation}
로 정의한다. 수렴한 HF solution에서는 $f=0$이다. Target energy 자체는 이 constraint에
대해 stationary하지 않으므로 adjoint multiplier $z$를 도입한 reduced Lagrangian을
\begin{equation}
  \mathcal L(C,z;\vR)
  =E_{\mathrm{DC-DFA}}(C;\vR)
   +2\operatorname{Re}\sum_{\vk ai}w_\vk
    z_{ai}^{\vk *}f_{ai}^{\vk}(C;\vR)
  \label{eq:dc-lagrangian}
\end{equation}
로 둔다. MO orthonormality multiplier는 아래의 energy-weighted density와 overlap/Pulay
항에 흡수하여 표기한다.

Orbital rotation에 대한 stationarity
\begin{equation}
  \frac{\partial\mathcal L}{\partial\kappa}=0
  \label{eq:lagrangian-orbital-stationarity}
\end{equation}
를 적용하면 식~\eqref{eq:zvector-equation}의 Z-vector equation을 얻는다. 따라서 total
derivative에서 명시적인 $\dd C/\dd R_{Ax}$ 항은 사라지고
\begin{equation}
  \frac{\dd E_{\mathrm{DC-DFA}}}{\dd R_{Ax}}
  =\left.\frac{\partial\mathcal L}{\partial R_{Ax}}\right|_{C,z}
  =g_{Ax}^{\mathrm{explicit}}+g_{Ax}^{\mathrm{response}}
  \label{eq:dc-gradient-master}
\end{equation}
만 남는다.

\subsection{Explicit target contribution}

Density matrix를 고정한 target energy의 explicit derivative를 다음처럼 분해한다.
\begin{equation}
  g_{Ax}^{\mathrm{explicit}}
  =g_{Ax}^{1e}+g_{Ax}^{\mathrm H}+g_{Ax}^{\mathrm{xc}}
   +g_{Ax}^{\mathrm{Pulay,t}}+g_{Ax}^{\mathrm{NN}}.
  \label{eq:explicit-gradient-decomposition}
\end{equation}

One-electron term은
\begin{equation}
  g_{Ax}^{1e}
  =\frac{1}{N_k}\sum_\vk
    \Tr[D^{\vk}h^{\vk,Ax}]
  \label{eq:one-electron-gradient}
\end{equation}
이다. $h^{Ax}$에는 kinetic, local GTH, nonlocal GTH의 AO matrix derivative가 모두
포함된다.

Hartree term은 density를 고정한 Coulomb matrix derivative로
\begin{equation}
  g_{Ax}^{\mathrm H}
  =\frac{1}{2N_k}\sum_\vk
    \Tr[D^{\vk}J^{\vk,Ax}[D]]
  \label{eq:hartree-gradient}
\end{equation}
이다. $J^{Ax}[D]$는 FFTDF에서는 AO/grid와 Coulomb-kernel derivative, GDF/RSDF에서는
three-center source와 two-center metric derivative를 포함한다.

Semilocal XC explicit derivative는 density matrix를 고정하고 AO와 quadrature representation을
미분한 값이다.
\begin{equation}
  g_{Ax}^{\mathrm{xc}}
  =\left.
    \frac{\partial E_{\mathrm{xc}}[D;\vR]}{\partial R_{Ax}}
   \right|_D.
  \label{eq:xc-explicit-gradient}
\end{equation}
Atom-centered Gaussian AO를 uniform periodic grid에서 평가하므로 AO center derivative와
periodic image/Bloch phase convention이 central value와 같아야 한다. SCAN/r2SCAN에서는
$\tau$ channel이 추가된다.

\subsection{Target overlap/Pulay term}

Target Fock의 occupied--occupied block으로 target energy-weighted density를
\begin{equation}
  W_{\mathrm t}^{\vk}
  =2C_o^{\vk}
    \left(C_o^{\vk\dagger}F_{\mathrm t}^{\vk}C_o^{\vk}\right)
    C_o^{\vk\dagger}
  \label{eq:target-energy-weighted-density}
\end{equation}
로 정의한다. 그러면 target overlap contribution은
\begin{equation}
  g_{Ax}^{\mathrm{Pulay,t}}
  =-\frac{1}{N_k}\sum_\vk
    \Tr[W_{\mathrm t}^{\vk}S^{\vk,Ax}]
  \label{eq:target-overlap-pulay}
\end{equation}
이다. AO-row derivative formulation에서는 식~\eqref{eq:one-electron-gradient}--
\eqref{eq:target-overlap-pulay}의 일부가 atom-row contraction으로 재배치될 수 있지만,
합계는 동일해야 한다.

Ion--ion 항은 GTH ionic charge와 periodic Ewald convention으로
\begin{equation}
  g_{Ax}^{\mathrm{NN}}
  =\frac{\partial E_{\mathrm{NN}}}{\partial R_{Ax}}
  \label{eq:nuclear-repulsion-gradient}
\end{equation}
이다.

\subsection{HF orbital-response contribution}

식~\eqref{eq:zvector-contraction}을 index로 쓰면
\begin{equation}
  g_{Ax}^{\mathrm{response}}
  =2\operatorname{Re}\sum_{\vk ai}w_\vk
   z_{ai}^{\vk *}B_{ai}^{\vk,Ax}.
  \label{eq:response-gradient}
\end{equation}
여기서 $B^{Ax}$는 source HF equation의 nuclear derivative다. Target PBE와 MGGA가
바뀌어도 $B^{Ax}$와 HF Hessian의 물리적 정의는 바뀌지 않는다. 바뀌는 것은 target
VO gradient $g$와 explicit XC derivative이므로 같은 response infrastructure를 재사용할
수 있다.

\subsection{Backend-independent assembly}

최종 package decomposition은
\begin{align}
  g_{Ax}^{\mathrm{elec,explicit}}
  &=g_{Ax}^{1e}+g_{Ax}^{\mathrm H}+g_{Ax}^{\mathrm{xc}}
    +g_{Ax}^{\mathrm{Pulay,t}},\\
  g_{Ax}
  &=g_{Ax}^{\mathrm{elec,explicit}}
    +g_{Ax}^{\mathrm{NN}}+g_{Ax}^{\mathrm{response}}
  \label{eq:final-gradient-assembly}
\end{align}
로 기록한다. FFTDF, GDF, RSDF의 차이는 $J/K$ value와 derivative provider에 있고,
Z-vector equation 및 final assembly contract는 공통이다.

\begin{remark}
Full-energy finite difference는 $\vR\pm h\mathbf e_{Ax}$ 각각에서 source HF를 다시
수렴시키고 target energy를 새로 평가해야 한다. Central density를 고정한 component FD는
식~\eqref{eq:xc-explicit-gradient} 같은 개별 kernel 진단에는 유용하지만 전체 analytic
gradient의 reference가 아니다.
\end{remark}

\subsection{Translational invariance}

고정된 lattice 안에서 모든 원자를 같은 벡터만큼 이동하면 energy는 변하지 않아야 한다.
따라서 converged representation에서
\begin{equation}
  \sum_A\mathbf g_A=\mathbf0
  \label{eq:translation-sum-rule}
\end{equation}
이다. 큰 net gradient는 AO image derivative, auxiliary-center response, XC grid aliasing,
$G=0$ 처리 또는 k-point normalization 누락을 찾는 강력한 diagnostic이다.
